\documentclass[10pt,a4paper]{article}
\usepackage[margin=21mm]{geometry}
\usepackage{amsmath,amssymb,booktabs,longtable}
\usepackage[colorlinks=true,urlcolor=blue,linkcolor=black]{hyperref}
\title{A064 v0.2.0 --- PKLSA-Initiated Euler Flow\\Separatrix Candidate and Material-Transport Falsifier}
\author{SST Workbench numerical research prototype}
\date{15 October 2026}
\begin{document}\maketitle
\begin{abstract}
An evolving, periodic, incompressible 3-D Euler numerical lane replaces the repeated frozen velocity fields in A064 v0.1.1. Original, unqualified 3-D KnotPlot/PKLSA centerline coordinates of $3_1$, $5_2$ and $6_1$ are used only as initial-condition shapes. The workflow independently discovers enstrophy-level candidate boundaries before auditing normal relative flux and tracer crossing. The findings concern the selected model-dependent level sets, not global invariant manifolds and not physical SST evidence. Framework v1.0.7 FREEZE/CERTIFY has not run.
\end{abstract}
\section{Source and evidence classification}
Input centerlines: three raw native knot polygons, retained with SHA-256 hashes. E010 v0.3.1/E011 v0.3.0 certify static geometry at their own admitted provider anchors; these bundled native polylines have \emph{not} been mapped to E011 STATIC\_READY anchor IDs. E011 additionally reports $3_1$ and $5_2$ cross-provider robust and $6_1$ cross-provider sensitive in static metrics. No physical dynamical velocity field has been imported.
\par All values in the numerical solver use \emph{dimensionless computational units}. The NPZ uses an arbitrary $L_0=1\,\mathrm m$, $T_0=1\,\mathrm s$ bookkeeping convention for compatibility; this is not a measurement nor calibration of $r_c$ or $\mathbf v_{\!\boldsymbol{\circlearrowleft}}$.
\section{Equations and numerical scheme}
For an inviscid incompressible Euler fluid without external forces:
\begin{align}
\nabla\cdot \mathbf u &= 0, &
\partial_t \mathbf u + (\mathbf u\cdot\nabla)\mathbf u &= -\frac1\rho\nabla p, \\
\boldsymbol\omega &= \nabla\times\mathbf u, &
\partial_t\mathbf u &= \mathsf P(\mathbf u\times\boldsymbol\omega),
\end{align}
where $\mathsf P$ is the Leray projection onto solenoidal fields. Spatial differentiation uses periodic 3-D Fourier modes with strict $2/3$ truncation; time stepping uses RK4 and a CFL limit of $0.16$. The initial knotted tube is a Gaussian-smoothed line-vorticity distribution projected into a divergence-free velocity field. It is a \emph{model assumption}, not a unique physical PKLSA reconstruction.
\par The solver monitors energy $E=\frac12\langle|\mathbf u|^2\rangle$, helicity $H=\langle\mathbf u\cdot\boldsymbol\omega\rangle$, divergence and
\begin{equation}
\frac{D|\boldsymbol\omega|^2}{Dt}=2\omega_i\omega_j\partial_j u_i\,,\quad \frac D{Dt}=\partial_t+\mathbf u\cdot\nabla.
\label{eq:stretch}
\end{equation}
For a discovered level set $F(\mathbf x,t)=|\boldsymbol\omega|^2-\lambda$, its relative normal speed is
\begin{equation}
V_n^{\rm rel}=\frac{\partial_t F+\mathbf u\cdot\nabla F}{|\nabla F|}.
\end{equation}
A vanishing value is necessary for this level set to be material; $\nabla\cdot\mathbf u=0$ alone does \emph{not} establish this. Equation (\ref{eq:stretch}) explains how vortex stretching can alter enstrophy level sets without destroying the material character of actual vortex tubes.
\section{Discovery and independent test}
\textbf{Phase D}: At $t=0$, choose the first closed, watertight $|\boldsymbol\omega|^2$ isosurface in the fixed percentile order $65,75,85,92$. Save surface vertices, triangle faces, candidate metadata and SHA-256 seals \emph{before} computing transport. No candidate tuning using tracer outcomes is allowed.
\par \textbf{Phase T}: Track paired inside/outside initial tracers via RK4 through the evolving velocity snapshots, then independently compare membership relative to the same Eulerian level $\lambda$ at the final time. Report relative normal speed, crossings and numerical refinement. These observables test an enstrophy feature; full three-dimensional stable/unstable manifold extraction is a distinct unresolved gate.
\section{Numerical outcomes (13 time frames; $N=24$)}
\begin{center}\small
\begin{tabular}{lrrrrr}\toprule
Knot & $\Delta E/E$ & $\max|\nabla\cdot u|$ & $Q_{95}/U$ & Crossings & Stretch res.\\\midrule
$3_1$ & 3.64e-10 & 2.07e-14 & 0.554 & 40/218 & 0.078 \\
\bottomrule\end{tabular}\end{center}

All three candidate enstrophy isosurfaces are nonmaterial within their respective \emph{periodic numerical models} at this screening resolution. This is not evidence against the existence of some \emph{other} invariant material surface.
\par The centered, three-snapshot vortex-stretching identity residuals are exploratory numerical diagnostics. The $6_1$ residual is larger and prevents strong qualification even though energy and divergence remain well controlled.
\section{Resolution and hard scientific gates}
For the $3_1$ solver, a two-grid ($N=24,32$) velocity comparison gave $L^2$ discrepancies of 1.17\% initially and 1.86\% finally. Two grids do not establish asymptotic convergence, and the fit between the core radius and cell width has not been independently calibrated.
\begin{itemize}
\item Software control suite: 17 passing tests, including the stationary ABC Beltrami Euler solution and RK4 refinement.
\item Time-varying NPZ generator, source hashes and candidate discovery seal: implemented.
\item Physical E011 provider-anchor propagation and native independent source agreement: \textbf{UNRESOLVED}.
\item Spatial/temporal convergence of transport flux and crossing fractions: \textbf{UNRESOLVED}.
\item Material versus dynamically invariant separatrix manifold: \textbf{UNRESOLVED}.
\item Framework v1.0.7 official FREEZE/CERTIFY and blind reveal ledger: \textbf{NOT EXECUTED}.
\end{itemize}
\section{Falsifiable next gates}
\begin{enumerate}
\item Replace raw centerlines by every eligible E011 provider anchor; preserve geometry hashes and independent-provider classes. Do not optimize carrier choice on a transport observable.
\item Run $N=24,32,48,64$ with controlled core width, and CFL/temporal ladders. Check normal flux and paired tracer statistics rather than energy drift alone.
\item Extract actual finite-time hyperbolic/elliptic transport structures or invariant manifolds in a suitable frame; validate material transport under time-varying rigid coordinate transformations.
\item Compare periodic Euler dynamics with an independently implemented, nonperiodic finite-core vortex evolution and with available externally generated snapshot data. No cross-model equivalence is presumed.
\end{enumerate}
\section{Reproduction}\texttt{run\_all.cmd EULER\_RUN}; \texttt{run\_all.cmd EULER\_BUILD --force}; \texttt{run\_all.cmd SELFTEST}. A model-NPZ source sidecar contains the input SHA-256, model parameters, timing, solver diagnostics and science-evidence class. Running the pipeline does not imply SST's physical constants are derived or validated.
\begin{thebibliography}{9}
\bibitem{helmholtz1858} H. von Helmholtz (1858), ``\"Uber Integrale der hydrodynamischen Gleichungen, welche den Wirbelbewegungen entsprechen'', \emph{Journal f\"ur die reine und angewandte Mathematik} \textbf{55}, 25--55. \url{https://doi.org/10.1515/crll.1858.55.25}.
\bibitem{saffman1992} P. G. Saffman (1992), \emph{Vortex Dynamics}, Cambridge University Press. \url{https://doi.org/10.1017/CBO9780511624063}.
\bibitem{hou2007} T. Y. Hou and R. Li (2007), ``Computing nearly singular solutions using pseudo-spectral methods'', \emph{Journal of Computational Physics} \textbf{226}, 379--397. \url{https://doi.org/10.1016/j.jcp.2007.04.015}.
\bibitem{haller2015} G. Haller (2015), ``Lagrangian Coherent Structures'', \emph{Annual Review of Fluid Mechanics} \textbf{47}, 137--162. \url{https://doi.org/10.1146/annurev-fluid-010313-141322}.
\end{thebibliography}
\end{document}
